\documentclass[11pt,aspectratio=169]{beamer}
\usetheme{metropolis}

\usepackage{fontspec}
\usepackage{booktabs}
\usepackage{graphicx}
\graphicspath{{imagens/}}

\newcommand{\placeholder}[2]{%
  \fbox{\parbox[c][#2][c]{#1}{\centering\textbf{PLACEHOLDER}}}%
}

\title{Escalando uma autenticação legada}
\date{\today}
\author{José Ricardo Alves Figueirôa}
\institute{Jusfy}

\begin{document}

\maketitle

\begin{frame}{whoami}
  \begin{columns}[T,onlytextwidth]
    \column{0.6\textwidth}
      \begin{itemize}
        \item José Ricardo Alves Figueirôa
        \item Techlead @ Jusfy
        \item Computer Engineer --- UFPE
        \item IoT Master --- Polytech Nice Côte d'Azur
      \end{itemize}
    \column{0.4\textwidth}
      
\includegraphics[width=\linewidth]{qrcode-border}
  \end{columns}
\end{frame}

\begin{frame}{Objetivo}
  \begin{itemize}
    \item Desacoplar autenticação do monolito
    \item Centralizar regras de negócio de login
    \item Remover boilerplate, duplicações de soluções e superfícies de ataque no login
  \end{itemize}
  \begin{center}
    \includegraphics[width=0.55\textwidth]{mesh-vs-star}
  \end{center}
\end{frame}

\section{Autenticação}

\begin{frame}{Autenticação (AuthN) vs Autorização (AuthZ)}
  \begin{center}
    \includegraphics[width=0.85\textwidth]{authz-vs-authn}
  \end{center}
\end{frame}

\begin{frame}{Três formas de provar quem você é}
  \begin{center}
    \includegraphics[width=0.8\textwidth]{auth-types}
  \end{center}
\end{frame}

\section{JWTs}

\begin{frame}{JWT}
  JWT (JSON Web Token) é um formato padronizado para transportar
  informações entre sistemas de forma compacta e verificável.
  \begin{center}
    \includegraphics[width=0.6\textwidth]{jwt}
  \end{center}
\end{frame}

\begin{frame}{Bearer Token}
  \begin{itemize}
    \item Bearer Authentication é um mecanismo de autenticação HTTP em
      que o cliente envia um token junto com a requisição.
    \item O nome ``Bearer'' significa literalmente: ``Quem portar esse
      token pode usá-lo.''
    \item A API não verifica quem está segurando o token; ela apenas
      verifica se o token é válido.
  \end{itemize}
\end{frame}

\section{Hydra}

\begin{frame}{O que é o Ory Hydra?}
  É o server de autorização OAuth2/OIDC, desenvolvido e mantido pela Ory.
  \begin{itemize}
    \item Valida tokens
    \item Emite tokens
    \item Assina tokens
    \item Revoga tokens
  \end{itemize}
\end{frame}

\begin{frame}{Topologia}
  \begin{center}
    \placeholder{0.9\textwidth}{5cm}
  \end{center}
\end{frame}

\section{OAuth2}

\begin{frame}{Client credentials}
  \begin{center}
    \includegraphics[width=0.85\textwidth]{m2m-flow}
  \end{center}
\end{frame}

\begin{frame}{PKCE}
  \begin{center}
    \includegraphics[width=0.7\textwidth]{h2m-flow}
  \end{center}
\end{frame}

\begin{frame}{Fluxo Implementado}
  \begin{center}
    \includegraphics[width=0.5\textwidth]{pkce-real}
  \end{center}
\end{frame}

\section{Arquitetura}

\begin{frame}{Arquitetura}
  \begin{center}
    \placeholder{0.9\textwidth}{5cm}
  \end{center}
\end{frame}

\begin{frame}{Arquitetura}
  \begin{center}
    \placeholder{0.9\textwidth}{5cm}
  \end{center}
\end{frame}

\begin{frame}[standout]
  Obrigado!

  \vspace{1em}
  
\includegraphics[width=0.2\textwidth]{qrcode-borderless}
\end{frame}

\end{document}
